\section{Introduction}
\label{sec:intro}

The Internet-of-Things (IoT) ecosystem comprises billions of deployed devices---routers, IP cameras, smart-home hubs, industrial controllers---most running embedded Linux with network-facing services that are routinely found vulnerable. CVE databases document tens of thousands of flaws in IoT firmware, from stack buffer overflows to authentication bypasses~\cite{cvedetails}. Reproducing a vulnerability---turning a CVE identifier into a controlled, observable trigger---is the prerequisite for patch validation, regression testing, and security research. Today this is a labor-intensive task demanding rare expertise spanning firmware reverse engineering, cross-architecture emulation, and binary exploitation.

LLM agents have shown strong results on software engineering~\cite{jimenez2024swebench}, code generation~\cite{chen2021codex}, and a rapidly growing set of cybersecurity tasks: vulnerability reproduction~\cite{wang2025cybergym}, exploitation~\cite{exploitgym,exploitbench}, kernel N-day reproduction~\cite{krepror}, and end-to-end patching~\cite{cybergyme2e}. A shared, often implicit, assumption underlies these benchmarks: the target is \emph{already executable}---provided as source code, a build script, or a sanitizer-instrumented Docker image. IoT firmware violates this assumption at every step. The artifact is a binary blob in a proprietary container format, compiled for a non-x86 instruction set, dependent on absent NVRAM state and hardware peripherals, and exposing the vulnerable service only after it is \emph{rehosted} under emulation. None of the capabilities required to cross this gap are exercised by existing benchmarks.

Starting from a CVE number and a device name, reproducing an IoT vulnerability demands six stages (\cref{sec:pipeline}): gather intelligence on the affected product and vulnerable endpoint; locate and download the matching firmware image; extract the filesystem from a proprietary, possibly encrypted, container; identify the architecture and the vulnerable binary; rehost the service under emulation; and trigger the flaw, optionally developing a full exploit and patch. Each stage hides distinct roadblocks---encrypted payloads, missing shared libraries, peripheral faults, NVRAM initialization---that determine whether reproduction is even possible. A single binary pass/fail verdict cannot reveal \emph{where} an agent fails, and therefore cannot guide progress.

We evaluate three frontier LLM agents over 372 runs spanning 124 IoT CVEs. Agents are competent at the early pipeline---fingerprinting firmware format (89\%) and architecture (82\%) and frequently downloading and extracting real firmware---but they systematically stall at rehosting. Only 12\% of runs invoke any emulator at all, full-system rehosting is attempted in just 2\% of runs, and successful rehosting is effectively absent. Instead, in 62\% of runs agents substitute a \emph{host-native simulation}: a freshly written program that mimics the vulnerable code path (e.g., an \texttt{strcpy} into a fixed buffer) and ``demonstrates'' a crash without ever executing the real firmware binary. This behavior inflates apparent success under naive metrics while leaving the actual vulnerability unreproduced. We call this systematic failure the \emph{rehosting gap} and argue it---not exploit development---is the dominant barrier to autonomous IoT security research.

\sys addresses this gap with four contributions. First, a pipeline-grounded roadblock taxonomy (\cref{sec:pipeline}) of 7 roadblocks, one per pipeline stage, each annotated with difficulty, state-of-the-art tooling, and a deterministic verification criterion. Second, the Roadblock Resolution Score (\cref{sec:rrs}), a normalized, dependency-aware metric that credits partial progress; we prove it is monotone with respect to the roadblock dependency graph, so it cannot reward an agent for skipping ahead. Third, the benchmark itself (\cref{sec:benchmark}): 124 real-world IoT CVEs across 8 device families and 4 architectures, with deterministic verification and no LLM-as-judge. Fourth, an empirical characterization of the rehosting gap (\cref{sec:experiments}) across 372 runs, showing that rehosting is the bottleneck, that simulation substitution is the dominant failure mode, and that current pass/fail metrics overstate true reproduction capability.
